\chapter{Introduction}\label{ch:intro}

\section{Motivation and gap}

Generative diffusion models dominate the synthesis of high-dimensional data
\citep{sohl2015deep, ho2020denoising, song2021scorebased}: a forward process
noises the data, a learned reverse process rebuilds it, and the one object
the reverse process needs is the \emph{score}
$\score(x,t) = \nabla_x \log p_t(x)$ of the noised law
\citep{anderson1982reverse}. In practice the score is a network trained by
denoising score matching \citep{vincent2011connection}, and for essentially
every data law of interest the score it should have produced is unavailable
\citep{kadkhodaie2024generalization}. Exactly solvable models supply such a
reference \citep{biroli2023generative, biroli2024dynamical}, but treat data
as unstructured samples, so the internal time of a \emph{sequence} does not
appear in them, and the classical inference literature on Gauss--Markov chains
\citep{kalman1960new, mezard2009information} does not ask how the exact
posterior, hence the exact score, deforms as the noise is swept from zero to
saturation. This thesis works in that gap, on sequences generated by a
first-order Markov chain and corrupted coordinatewise by the
Ornstein--Uhlenbeck channel. The corruption adds no edges among the latent
variables, so the posterior factor graph is a tree on which sum--product is
exact, and the \emph{joint} score of the whole noisy sequence, the object
one-frame scores are blind to (Chapter~\ref{ch:gaussian}), is an exact
recursion even for non-Gaussian innovations. The thesis computes it exactly
where a closed form exists, approximates it in a measured way where none
does, and learns it from noised data when the dynamics are unknown.

\section{Research questions}\label{sec:rq}

\begin{itemize}[nosep, leftmargin=*]
  \item[\textbf{RQ1.}] For a stationary Gaussian AR(1) chain under
  coordinatewise OU noise, what is the joint score, and how do its couplings
  vary with diffusion time?
  \item[\textbf{RQ2.}] How is the same score computed by belief propagation,
  and at what cost?
  \item[\textbf{RQ3.}] For non-Gaussian innovations, how accurate is a
  Gaussian-message approximation against a validated numerical reference?
  \item[\textbf{RQ4.}] How fast does the influence of remote observations on
  the score decay, and what does this say about local approximations?
  \item[\textbf{RQ5.}] Under what assumptions can an \emph{unknown} transition
  law be recovered from noised sequences alone, and which checks make that
  reliable?
  \item[\textbf{RQ6.}] Where the Markov assumption holds exactly, what is
  supplying it worth against networks trained on the same data and given
  none of it?
\end{itemize}

\section{Contributions and structure}

\begin{enumerate}[nosep, leftmargin=*]
  \item \textbf{The Gaussian case worked out completely:} closed-form score,
  spectrum and coupling lifecycle, chain belief propagation identified with
  Kalman filtering and RTS smoothing, and the measured price of ignoring
  distant evidence (Chapters~\ref{ch:gaussian}--\ref{ch:gaussian-bp}; RQ1,
  RQ2, RQ4).
  \item \textbf{What Gaussian closure computes off the Gaussian case:} the
  LMMSE estimator of the covariance-matched model
  (Theorem~\ref{thm:l-lmmse}), which makes the measured closure gap the price
  of discarding structure beyond second moments (Chapter~\ref{ch:laplace};
  RQ3).
  \item \textbf{Learning the chain, and what structure is worth:} an unknown
  kernel fitted by EM with Fisher's identity supplying the exact gradient
  from one forward--backward pass, its two convergence rates, and a measured
  value of $\ratiolo$--$\ratiohi\times$ against an unstructured network and
  $\structratiolo$--$\structratiohi\times$ against one given locality, not a
  ranking of algorithms (Chapters~\ref{ch:em-method}--\ref{ch:em-results};
  RQ5, RQ6).
\end{enumerate}
Chapter~\ref{ch:statmech} reviews the literature; Chapter~\ref{ch:conclusions}
answers the six questions.
